\vsssub
\subsubsection{~网格} \label{sec:grids}
\vsssub

% fig_grids_1

\input{zh/sys/fig_grd1_zh}

为便于编程并节省开销，空间网格和谱网格分开考虑。这种做法受
第~\ref{chapt:num} 章所述分裂技术启发。对于空间传播，使用一个简单的
“矩形”空间网格，如图~\ref{fig:grids_1} 所示。该网格可以是笛卡尔
``$(x,y)$'' 网格、球面网格（经纬度上为规则步长）、曲线网格，或基于三角形
的网格。在球面网格中，程序中始终用计数器 {\F ix} 表示经度，用计数器
{\F iy} 表示纬度，并采用相应网格维度 {\F (nx,ny)}。所有空间场数组都在
代码中动态分配；相应工作数组通常为自动数组，以允许线程安全代码。全局
应用中网格闭合由模式内部处理，不需要用户干预。为简化导数计算，特别是
流的导数计算，外侧网格点（{\F ix=1,nx}，除非网格为全局网格）以及
（{\F iy=1,ny}）会被视为陆点、非活动点或活动边界点。因此，除基于三角形
的网格外，最小网格尺寸为 {\F nx=3, ny=3}。在后一种情况下，所有节点作为
维度为 {\F nx} 的长向量列出，同时令 {\F ny=1}，从而保持相同代码结构。
输入数组通常假定具有如下形式

\vspace{\baselineskip} \centerline{\F array(nx,ny) ,} \vspace{\baselineskip}

\noindent
并逐行读取（另见第~\ref{chapt:run} 章）。不过，在程序内部，它们通常以
旋转后的索引存储：

\vspace{\baselineskip} \centerline{\F array(ny,nx) .} \vspace{\baselineskip}

\noindent
这样更容易提供全局闭合；全局闭合通常需要扩展 $x$ 轴。此外，这类二维数组
通常按一维数组处理，以增加向量长度。数组 {\F array}、其等价一维形式
{\F varray} 以及 {\F ixy} 定义为

\vspace{\baselineskip}
\centerline{\F array(my,mx) , varray(my*mx) ,}
\centerline{\F ixy = iy + (ix-1)*my .}
\vspace{\baselineskip}

\noindent
注意，网格的这种表示只在模式{\it 内部}使用。

% fig:grids_2

\input{zh/sys/fig_grd2_zh}

给定空间网格点 {\F (ix,iy)} 处的谱网格以类似方式定义，使用方向计数器
{\F ith} 和波数计数器 {\F ik}（图~\ref{fig:grids_2}）。谱网格大小通过
动态分配设置。与空间网格一样，谱 {\F a} 的内部描述定义为

\vspace{\baselineskip}
\centerline{\F a(nth,nk) ,}
\vspace{\baselineskip}

\noindent
程序中通篇使用等价的一维数组。在模式内部，方向始终为笛卡尔方向，
$\theta = 0^\circ$ 对应从西向东传播（正 $x$ 或 {\F ix} 方向），
$\theta = 90^\circ$ 对应从南向北传播（正 $y$ 或 {\F iy} 方向）。输出方向
使用其他约定，如第~\ref{chapt:run} 章所述。

波谱存储占据模式所需内存的大部分，因为所采用的分裂技术保证模式任一部分
只在整个波场的一个小子集上运行。为尽量减少所需内存，只存储实际海点处的
谱。这里的海点定义为可能需要谱的点。这包括活动边界点和被冰覆盖的海点。
出于归档目的，使用计数器 {\F isea} 定义一维海点网格。随后谱存储为

\vspace{\baselineskip}
\centerline{\F a(ith,ik,isea) .}
\vspace{\baselineskip}

\noindent
图~\ref{fig:grids_3} 给出了该存储网格相对于图~\ref{fig:grids_1} 中完整
网格布局的示例。显然，存储网格和完整空间网格之间的关系需要一些簿记。
为此定义了两个“映射”{\F mapfs} 和 {\F mapsf}。

\vspace{\baselineskip}
\centerline{\F mapsf(isea,1) = ix ,}
\centerline{\F mapsf(isea,2) = iy ,}
\centerline{\F mapsf(isea,3) = ixy ,}
\centerline{\F mapfs(iy,ix) = vmapfs(ixy) = isea ,}
\vspace{\baselineskip}

% fig:grids_3

\input{zh/sys/fig_grd3_zh}

\noindent
其中，对于陆点，{\F mapfs(iy,ix) = 0}。最后，维护状态映射
{\F mapsta(iy,ix)} 和 {\F mapst2(iy,ix)}，用于识别海点、陆点、活动边界点
和冰点。{\F mapsta} 表示网格的主状态映射：

\vspace{\baselineskip} \noindent
\strut \hspace{25mm} {\F mapsta(iy,ix) = 0} \hspace{10mm} 表示排除点，\\
\strut \hspace{25mm} {\F mapsta(iy,ix) = 1} \hspace{10mm} 表示海点，\\
\strut \hspace{25mm} {\F mapsta(iy,ix) = 2} \hspace{10mm} 表示活动边界点。

\vspace{\baselineskip}

由于存在海冰而未纳入波浪模式计算的海点和活动边界点，会以相应的负状态
指示符（-1 或 -2）标记。{\F mapst2} 包含二级信息。对于排除点
{\F mapsta(iy,ix) = 0}，该映射区分陆点 {\F mapst2(iy,ix) = 0} 和其他排除点
{\F mapst2(iy,ix) = 1}。对于被禁用的海点 {\F mapsta(iy,ix) < 0}，
{\F mapst2} 中的连续位标识停用原因（位值 1 表示停用）。

\begin{center} \begin{tabular}{cl}
 bit & 标识 \\ \hline
  1  & 海冰覆盖     \\
  2  & 点已干出  \\
  3  & 移动网格中的陆地或嵌套中推断出的陆地 \\
  4  & 在双向嵌套中被掩膜
\end{tabular} \end{center}

这里还考虑了另外两点。首先，两个状态映射可以折叠为单个映射以便存储。
为保证这种存储与前一模式版本向后兼容，两个映射合并为单个映射 {\F maptmp}：

\vspace{\baselineskip}
\centerline{\F maptmp = mapsta + 8 * mapst2}
\vspace{\baselineskip}

\noindent
这是因为只有 {\F mapsta} 的前几位包含数据。NetCDF 文件中保存的正是该
映射 MAPTMP。原始映射可恢复为

\vspace{\baselineskip}
\centerline{\F mapsta = mod ( maptmp + 2 , 8 ) - 2}
\centerline{\F mapst2 = maptmp - mapsta}
\vspace{\baselineskip}

\noindent
其次，图形输出程序使用单个映射，以简化网格点状态的绘制。在图形文件中，
该映射定义为

\begin{center} \begin{tabular}{cl}
{map} & 表示 \\ \hline
  2 & 活动边界点 \\
  1 & 活动海点      \\
  0 & 陆点（包括在 {\F MAPST2} 中识别出的陆点） \\
 -1 & 被冰覆盖但为湿点的点 \\
 -2 & 未被冰覆盖的干点 \\
 -3 & 被冰覆盖的干点 \\
 -4 & 在双向嵌套方案中被掩膜的点 \\
 -5 & 其他禁用点
\end{tabular} \end{center}

\noindent
类似地，在网格准备程序 {\file ww3\_grid} 中，也可使用单个映射来简化处理。
该映射按如下方式区分各点：

\begin{center} \begin{tabular}{cl}
{map} & 表示 \\ \hline
  3 & 排除点 \\
  2 & 活动边界点 \\
  1 & 活动海点      \\
  0 & 陆点 
\end{tabular} \end{center}
